\chapter{Random Walks}

\begin{center}\itshape
A drunk man will find his way home, but a drunk bird may get lost forever.\\
\upshape(attributed to Shizuo Kakutani)
\end{center}

Toss a fair coin again and again, and move one step up for heads and one step down for tails. This is the simple
random walk, the most basic random process of all, and it already holds surprises. The walk returns to its start
with probability one, yet the expected time until it does is infinite. In a long game between two equally lucky
players, one of them is usually ahead for most of the time, and the last tie is more likely to be near the very
beginning or the very end than in the middle. In the plane the walk still comes home; in space it does not.

The tools are the reflection principle of Chapter~1, a little counting, and independence. With them we compute the
law of the first return and of the maximum, prove the arcsine law for the last tie, solve the gambler's ruin with
its durations, and prove Pólya's theorem on recurrence in every dimension. Then the walk turns into a tool: its
exit point solves the discrete Dirichlet problem, the ancestor of the link between Brownian motion and harmonic
functions, and a walk on assignments solves the satisfiability problem 2-SAT in quadratic time.

\begin{think}
A fair coin is tossed once a second for a year, and we follow the difference between heads and tails. How likely
is it that the last tie happened during the first week, so that one side has been ahead without interruption for
the remaining fifty-one weeks? Guess an order of magnitude before reading Section~\ref{sec:arcsine}.
\end{think}

\section{The simple random walk}

Let $X_1,X_2,\ldots$ be independent with $P(X_i=1)=P(X_i=-1)=\frac12$ (such a sequence exists by the admitted fact (A1)
of Chapter~5), and let $S_0=0$, $S_n=X_1+\dots+X_n$. An event that depends only on $X_1,\dots,X_n$ is described by a set of
paths $(s_0,\dots,s_n)$ with steps $\pm1$, and its probability is the number of these paths divided by $2^n$. A path with
$u$ up-steps and $d$ down-steps ends at $u-d$, so
\[
P(S_n=k)=\binom n{\frac{n+k}2}2^{-n}
\]
if $n+k$ is even and $|k|\le n$, and $0$ otherwise. In particular the walk can be at $0$ only at even times. Write
\[
u_{2n}=P(S_{2n}=0)=\binom{2n}n4^{-n},
\]
so that $u_0=1$, $u_2=\frac12$, $u_4=\frac38$. Stirling's formula (Tome~0, Theorems~5.6 and~6.4) gives $u_{2n}\sim\frac1{\sqrt{\pi n}}$, but the following
elementary bounds are all we shall need.

\begin{lemma}\label{lem:central}
For $n\ge1$, $\dfrac1{2\sqrt n}\le u_{2n}\le\dfrac1{\sqrt{2n+1}}$. More generally, for every $m\ge0$ and every $k$, $\binom mk2^{-m}\le\frac1{\sqrt{m+1}}$.
\end{lemma}
\begin{proof}
Since $\binom{2n}n=\frac{(2n)!}{n!\,n!}$ and $(2n)!=2^nn!\cdot1\cdot3\cdots(2n-1)$,
\[
u_{2n}=\prod_{j=1}^n\frac{2j-1}{2j}.
\]
As $\frac{2j-1}{2j}\le\frac{2j}{2j+1}$, we get $u_{2n}^2\le\prod_j\frac{2j-1}{2j}\cdot\frac{2j}{2j+1}=\prod_j\frac{2j-1}{2j+1}=\frac1{2n+1}$, the product telescoping. For $j\ge2$,
$\frac{2j-1}{2j}\ge\frac{2j-2}{2j-1}$, so
\[
\begin{aligned}u_{2n}^2&=\frac14\prod_{j=2}^n\Bigl(\frac{2j-1}{2j}\Bigr)^2\eqbr\ge\frac14\prod_{j=2}^n\frac{2j-2}{2j-1}\cdot\frac{2j-1}{2j}\\&=\frac14\prod_{j=2}^n\frac{j-1}j=\frac1{4n}.\end{aligned}
\]
For the general bound, $\binom m{k+1}/\binom mk=\frac{m-k}{k+1}$ is at least $1$ exactly when $k\le\frac{m-1}2$, so $\binom mk$ is largest at
$k=\lfloor\frac m2\rfloor$. If $m=2l$ the maximum is $u_{2l}\le\frac1{\sqrt{2l+1}}=\frac1{\sqrt{m+1}}$. If $m=2l+1$, then $\binom{2l+1}l=\frac12\binom{2l+2}{l+1}$, so the
maximum is $u_{2l+2}\le\frac1{\sqrt{2l+3}}\le\frac1{\sqrt{m+1}}$.
\end{proof}

The walk \emph{starts afresh} after any fixed time $k$: the increments $(S_{k+j}-S_k)_{j\ge0}=(X_{k+1}+\dots+X_{k+j})_{j\ge0}$ form a
simple random walk, independent of $X_1,\dots,X_k$. We shall use this at random times $T$ by splitting according to the value of
$T$, as in the next lemma, where $T=\min\{n\ge1:S_n=0\}$ is the time of the first return to $0$ ($T=\infty$ if there is none).
Write $f_{2n}=P(T=2n)$.

\begin{lemma}[first-return decomposition]\label{lem:renewal}
For $n\ge1$, $u_{2n}=\sum_{r=1}^nf_{2r}\,u_{2n-2r}$.
\end{lemma}
\begin{proof}
If $S_{2n}=0$ the walk has returned to $0$ by time $2n$, so $\{S_{2n}=0\}$ is the disjoint union of the events
$\{T=2r,\ S_{2n}=0\}$ for $1\le r\le n$. The event $\{T=2r\}$ depends only on $X_1,\dots,X_{2r}$, and on it $S_{2n}=S_{2n}-S_{2r}$, which
depends only on $X_{2r+1},\dots,X_{2n}$ and has the law of $S_{2n-2r}$. By independence, $P(T=2r,\ S_{2n}=0)=f_{2r}u_{2n-2r}$.
\end{proof}

\section{Reflection, first returns and the maximum}

The reflection principle (Lemma~\ref{lem:reflection}) counts the paths that touch $0$. Here is its main consequence for
the walk, an identity that looks like a coincidence.

\begin{theorem}[the main lemma]\label{thm:mainlemma}
For every $n\ge1$,
\begin{enumerate}
\item[(a)] $P(S_1\ne0,\ S_2\ne0,\ \dots,\ S_{2n}\ne0)=u_{2n}$;
\item[(b)] $P(S_1\ge0,\ S_2\ge0,\ \dots,\ S_{2n}\ge0)=u_{2n}$;
\item[(c)] $P(T>2n)=u_{2n}$, and $f_{2n}=u_{2n-2}-u_{2n}=\dfrac{u_{2n}}{2n-1}$.
\end{enumerate}
\end{theorem}
\begin{proof}
(a) By symmetry, the probability is $2P(S_1>0,\dots,S_{2n}>0)$. A path with $S_1>0,\dots,S_{2n}>0$ starts with an up-step and
ends at some $2k$ with $1\le k\le n$; after the first step it is a path from $(1,1)$ to $(2n,2k)$ that does not touch $0$.
There are $N(1,2k)$ paths of length $2n-1$ from $1$ to $2k$, where $N(x,y)=\binom{2n-1}{(2n-1+y-x)/2}$, and by the reflection
principle $N(-1,2k)$ of them touch $0$. So the number of paths with $S_1>0,\dots,S_{2n}>0$ is
\[
\sum_{k=1}^n\Bigl[\binom{2n-1}{n+k-1}-\binom{2n-1}{n+k}\Bigr]=\binom{2n-1}n,
\]
the sum telescoping (the last term $\binom{2n-1}{2n}$ is $0$). Since $\binom{2n-1}n=\frac12\binom{2n}n$, the probability of (a) is $2\cdot\frac12\binom{2n}n2^{-2n}=u_{2n}$.

(b) A path with $S_1>0,\dots,S_{2n}>0$ is an up-step followed by a path $s'$ of length $2n-1$ from $1$ that stays $\ge1$,
that is, by $1+s''$ where $s''$ is a path of length $2n-1$ from $0$ with $s''_j\ge0$ for all $j$. So
\[
\begin{aligned}\tfrac12u_{2n}&=P(S_1>0,\dots,S_{2n}>0)\eqbr=\tfrac12P(S_1\ge0,\dots,S_{2n-1}\ge0).\end{aligned}
\]
Since $2n-1$ is odd, $S_{2n-1}\ge0$ forces $S_{2n-1}\ge1$ and then $S_{2n}\ge0$; so the last probability equals $P(S_1\ge0,\dots,S_{2n}\ge0)$.

(c) $\{T>2n\}$ is the event of (a). As $T$ is even, $f_{2n}=P(T>2n-2)-P(T>2n)=u_{2n-2}-u_{2n}$, and $u_{2n}=u_{2n-2}\cdot\frac{2n-1}{2n}$
gives $u_{2n-2}-u_{2n}=u_{2n}\bigl(\frac{2n}{2n-1}-1\bigr)=\frac{u_{2n}}{2n-1}$.
\end{proof}

\begin{corollary}\label{cor:return1d}
The simple random walk on $\Z$ returns to $0$ with probability $1$, but $\mathbb E[T]=\infty$.
\end{corollary}
\begin{proof}
$P(T=\infty)\le P(T>2n)=u_{2n}\le\frac1{\sqrt{2n+1}}\to0$. And $\mathbb E[T]=\sum_{n\ge0}P(T>n)\ge\sum_nP(T>2n)=\sum_nu_{2n}\ge\sum_{n\ge1}\frac1{2\sqrt n}=\infty$.
\end{proof}

So a fair game between two players always comes back to a tie, but you should not wait for it: the expected waiting
time is infinite, and $P(T>2n)\approx\frac1{\sqrt{\pi n}}$ decays very slowly. After a million tosses without a tie, the
chance of a further million without one is still $\frac{u_{2\cdot10^6}}{u_{10^6}}\approx\frac1{\sqrt2}$.

The same trick of reflecting a path after it first reaches a level gives the law of the maximum.

\begin{proposition}[the maximum]\label{prop:maximum}
Let $M_n=\max(S_0,\dots,S_n)$. For integers $r\ge1$ and $y\le r$, $P(M_n\ge r,\ S_n=y)=P(S_n=2r-y)$. Consequently
\[
\begin{gathered}P(M_n\ge r)=P(S_n\ge r)+P(S_n>r),\qq P(M_n=r)=P(S_n=r)+P(S_n=r+1).\end{gathered}
\]
\end{proposition}
\begin{proof}
For a path with $M_n\ge r$ and $s_n=y$, let $\tau$ be the first time it reaches $r$, and reflect the part after $\tau$ in the
line at height $r$: $s'_k=s_k$ for $k\le\tau$ and $s'_k=2r-s_k$ for $k\ge\tau$. This gives a path ending at $2r-y\ge r$. Conversely, a
path ending at $2r-y\ge r$ reaches $r$, and reflecting it after its first visit to $r$ gives back the original path,
because the first visit to $r$ happens at the same time $\tau$ for $s$ and $s'$. So the two sets of paths have the same size.
Summing over $y<r$, and adding the paths with $S_n\ge r$ (which all have $M_n\ge r$),
\[
\begin{aligned}P(M_n\ge r)&=P(S_n\ge r)+\sum_{y<r}P(S_n=2r-y)\eqbr=P(S_n\ge r)+P(S_n>r).\end{aligned}
\]
Subtracting the same identity for $r+1$ gives $P(M_n=r)$. For $r=0$ the formulas also hold, as both sides of the
first are $1$ and $P(M_n=0)=1-P(M_n\ge1)$.
\end{proof}

\begin{figure}[ht]
\centering
\fig{ch07-reflect}
\caption{Reflection at level $r=4$. After its first visit to $r$ at time $\tau$, the path (orange) is reflected in the
dashed line (blue). Paths that reach $r$ and end at $y$ correspond exactly to paths that end at $2r-y$.}
\label{fig:reflect}
\end{figure}

Exercise~\ref{ex:maxlaw} deduces from the central limit theorem that $\frac{M_n}{\sqrt n}$ converges in law to $|N|$, where $N$ is
$N(0,1)$. Chapter~11 proves the same statement for the maximum of Brownian motion.

\section{The arcsine law}\label{sec:arcsine}

Let $L_{2n}=\max\{k\le2n:S_k=0\}$ be the time of the last tie up to time $2n$.

\begin{theorem}[discrete arcsine law]\label{thm:arcsine}
For $0\le k\le n$, $P(L_{2n}=2k)=u_{2k}\,u_{2n-2k}$. In particular the law of $L_{2n}$ is symmetric about $n$.
\end{theorem}
\begin{proof}
$L_{2n}=2k$ means that $S_{2k}=0$ and that $S_j-S_{2k}\ne0$ for $2k<j\le2n$. The first event depends on $X_1,\dots,X_{2k}$; the second
is the event of Theorem~\ref{thm:mainlemma}(a) for the walk $(S_{2k+j}-S_{2k})_j$, which is independent of the first and has probability
$u_{2n-2k}$. The formula is unchanged when $k$ is replaced by $n-k$.
\end{proof}

The values $u_{2k}u_{2n-2k}$ are largest at $k=0$ and $k=n$ and smallest in the middle: the last tie is most likely to be
at the very beginning or the very end. In the limit the law has a density that explodes at both ends.

\begin{theorem}[arcsine law]\label{thm:arcsinelimit}
For $0<a<b<1$,
\[
\begin{aligned}P\Bigl(a\le\frac{L_{2n}}{2n}\le b\Bigr)&\xrightarrow[n\to\infty]{}\int_a^b\frac{\dd x}{\pi\sqrt{x(1-x)}}\\&=\frac2\pi\bigl(\arcsin\sqrt b-\arcsin\sqrt a\bigr).\end{aligned}
\]
\end{theorem}
\begin{proof}
Write $u_{2k}=\frac{1+\varepsilon_k}{\sqrt{\pi k}}$ with $\varepsilon_k\to0$ (Stirling). The event is $\{an\le k\le bn\}$ with $L_{2n}=2k$, and for such $k$ both $k\ge an$
and $n-k\ge(1-b)n$ tend to infinity, so
\[
\begin{gathered}u_{2k}u_{2n-2k}=\frac{1+\eta_{n,k}}{\pi\sqrt{k(n-k)}}=\frac{1+\eta_{n,k}}n\,g\Bigl(\frac kn\Bigr),\qq g(x)=\frac1{\pi\sqrt{x(1-x)}},\end{gathered}
\]
with $\max_{an\le k\le bn}|\eta_{n,k}|\to0$. Hence $P(a\le\frac{L_{2n}}{2n}\le b)$ differs from $\frac1n\sum_{an\le k\le bn}g(\frac kn)$ by at most $\max_k|\eta_{n,k}|$ times
this sum, which is bounded. The sum is a Riemann sum of the continuous function $g$ on $[a,b]$ (up to at most two
terms of size $\frac{\max g}n$ at the ends), and tends to $\int_a^bg$. Finally, with $x=\sin^2\theta$, $\dd x=2\sin\theta\cos\theta\dd\theta$ and
$\sqrt{x(1-x)}=\sin\theta\cos\theta$, so the integral is $\frac2\pi$ times the difference of the values of $\theta=\arcsin\sqrt x$.
\end{proof}

\begin{figure}[ht]
\centering
\fig{ch07-arcsine}
\caption{The law of the last tie $L_{2n}$ for $2n=100$ tosses: $P(L_{2n}=2k)=u_{2k}u_{2n-2k}$ (bars), against the arcsine density
$\frac1{\pi\sqrt{x(1-x)}}$ rescaled to the same axis (curve). The least likely place for the last tie is the middle.}
\label{fig:arcsine}
\end{figure}

For the coin tossed every second for a year, the first week is $\frac7{365}\approx0.02$ of the time, and the probability that
the last tie happened in it is about $\frac2\pi\arcsin\sqrt{0.02}\approx0.09$: about one chance in eleven that one player has led
without a break for the last fifty-one weeks. The probability that the last tie falls in the first quarter of the
year is $\frac2\pi\arcsin\frac12=\frac13$.

Exercise~\ref{ex:timepositive} proves that the number of steps during which the walk is above $0$ has the same law
$u_{2k}u_{2n-2k}$. So the most likely outcome of a long fair game is that one player is ahead almost all the time.

\begin{beauty}
Two different questions, ``when was the last tie?'' and ``how long was the walk positive?'', have the same answer,
the same discrete law $u_{2k}u_{2n-2k}$, and in the limit the same curve $\frac1{\pi\sqrt{x(1-x)}}$, the density of $\sin^2$ of a uniform angle.
Paul Lévy found the arcsine law for Brownian motion in 1939, and Erdős and Kac proved it for general sums in
1947; the walk shows it with nothing more than counting paths.
\end{beauty}

\section{The gambler's ruin}

Now let the walk be biased: $P(X_i=1)=p$ and $P(X_i=-1)=q=1-p$, with $0<p<1$. A gambler starts with $x$ euros, bets one
euro at a time, wins each bet with probability $p$, and stops when the fortune $x+S_n$ reaches $0$ (ruin) or $N$ (the
target), where $0\le x\le N$. Let
\[
\tau=\min\{n\ge0:x+S_n\in\{0,N\}\}.
\]
We write $P_x$ and $\mathbb E_x$ to recall the starting fortune.

\begin{lemma}\label{lem:exit}
$P_x(\tau>kN)\le(1-p^N)^k$ for all $k$. Hence $\tau<\infty$ almost surely and $\mathbb E_x[\tau]<\infty$.
\end{lemma}
\begin{proof}
Cut time into blocks of $N$ steps. If the $N$ steps of a block are all up-steps, the game has ended by the end of that
block, since from any fortune in $[0,N]$ it reaches $N$ within $N$ up-steps. So $\{\tau>kN\}$ is contained in the event that none of
the first $k$ blocks consists only of up-steps; the blocks are independent, each all-up with probability $p^N$. Then
$\mathbb E_x[\tau]=\sum_{m\ge0}P_x(\tau>m)\le N\sum_{k\ge0}(1-p^N)^k=Np^{-N}$.
\end{proof}

\begin{theorem}[gambler's ruin]\label{thm:ruin}
Let $r=\frac qp$. The probability $h(x)=P_x(x+S_\tau=N)$ of reaching the target and the expected duration $m(x)=\mathbb E_x[\tau]$ are
\[
\begin{gathered}h(x)=\begin{cases}\dfrac{r^x-1}{r^N-1}&(p\ne q),\\[2ex]\dfrac xN&(p=q),\end{cases}\qq m(x)=\begin{cases}\dfrac{x-Nh(x)}{q-p}&(p\ne q),\\[2ex]x(N-x)&(p=q).\end{cases}\end{gathered}
\]
\end{theorem}
\begin{proof}
\emph{First-step analysis.} The event $\{x+S_\tau=N\}$ is $\{(X_1,X_2,\ldots)\in A_x\}$ for a set $A_x$ of sequences of signs. For $0<x<N$ the game
does not end at time $0$, and $(X_1,X_2,\ldots)\in A_x$ if and only if either $X_1=1$ and $(X_2,X_3,\ldots)\in A_{x+1}$, or $X_1=-1$ and
$(X_2,X_3,\ldots)\in A_{x-1}$. The sequence $(X_2,X_3,\ldots)$ is independent of $X_1$ and has the same law as $(X_1,X_2,\ldots)$ (the two laws agree
on events fixing finitely many terms, a $\pi$-system, so they agree everywhere by Theorem~\ref{thm:uniqueness}). Hence
\[
\begin{gathered}h(x)=p\,h(x+1)+q\,h(x-1)\quad(0<x<N),\qq h(0)=0,\quad h(N)=1.\end{gathered}
\]
In the same way $\tau=1+\tau'$, where $\tau'$ is the duration of the game played with $(X_2,X_3,\ldots)$ from $x\pm1$, so for $0<x<N$
\[
\begin{gathered}m(x)=1+p\,m(x+1)+q\,m(x-1),\qq m(0)=m(N)=0,\end{gathered}
\]
all quantities being finite by Lemma~\ref{lem:exit}.

\emph{Solving.} Writing $1=p+q$, the first equation reads $p\bigl(h(x+1)-h(x)\bigr)=q\bigl(h(x)-h(x-1)\bigr)$, so the increments are
$h(x)-h(x-1)=r^{x-1}c$ for a constant $c$, and $h(x)=c(1+r+\dots+r^{x-1})$. The condition $h(N)=1$ fixes $c$, which gives the formula. For
$m$, one checks directly that the stated functions satisfy the equations: for $p=q$,
\[
\begin{aligned}x(N-x)&=1+\tfrac12(x+1)(N-x-1)\eqbr+\tfrac12(x-1)(N-x+1),\end{aligned}
\]
and for $p\ne q$, $\frac x{q-p}$ satisfies $1+\frac{p(x+1)+q(x-1)}{q-p}=\frac x{q-p}$ while $h$ satisfies the equation without the $1$, so their
combination does, with the right boundary values. The solution is unique: the difference $\delta$ of two solutions
satisfies $\delta(x)=p\delta(x+1)+q\delta(x-1)$ with $\delta(0)=\delta(N)=0$, and the computation just made for $h$ shows $\delta(x)=c(1+\dots+r^{x-1})$, so $c=0$.
\end{proof}

\begin{corollary}\label{cor:infinitebank}
For a gambler playing against an infinitely rich bank, the probability of eventual ruin from fortune $x\ge1$ is $1$
if $p\le q$ and $\bigl(\frac qp\bigr)^x$ if $p>q$.
\end{corollary}
\begin{proof}
If the walk from $x$ ever hits $0$, say at time $t$, then it hits $0$ before $N$ for every $N>x+t$. So the event of eventual ruin is
the increasing union over $N$ of $\{$hit $0$ before $N\}$, and its probability is $\lim_N\bigl(1-h_N(x)\bigr)$. When $p>q$, $r<1$ and $h_N(x)\to1-r^x$; when
$p<q$, $r>1$ and $h_N(x)\to0$; when $p=q$, $h_N(x)=\frac xN\to0$.
\end{proof}

This agrees with Exercise~\ref{ex:firstpassage}, which found the same numbers with generating functions.

\begin{example}
On a European roulette wheel a bet on red wins with probability $p=\frac{18}{37}$. A player with $100$ euros who bets one
euro at a time until reaching $200$ or ruin succeeds with probability $\frac{r^{100}-1}{r^{200}-1}\approx0.0045$, where $r=\frac{19}{18}$, after an
average of $3667$ bets. Betting $10$ euros at a time (a game from $10$ to $20$ in units of $10$) raises the chance to $0.37$,
and a single bet of $100$ euros to $\frac{18}{37}\approx0.49$. In an unfavourable game, timidity is fatal: every bet pays the house
its small edge, and many small bets let the law of large numbers work for the casino (Exercise~\ref{ex:unfair}).
\end{example}

\begin{think}
In a fair game ($p=\frac12$) from $x=100$ with target $200$, the success probability is $\frac12$ whatever the stake. What does the
stake change?
\end{think}

\section{Recurrence and Pólya's theorem}

The \emph{simple random walk on $\Z^d$} moves at each step to one of the $2d$ neighbours of its position, chosen uniformly:
$S_n=X_1+\dots+X_n$ with $X_i$ independent and uniform on $\{\pm e_1,\dots,\pm e_d\}$, where $e_1,\dots,e_d$ is the standard basis. Let $u_n=P(S_n=0)$,
let $T=\min\{n\ge1:S_n=0\}$, $F=P(T<\infty)$ the probability of ever returning, and $V=\#\{n\ge1:S_n=0\}$ the number of
returns. The walk is \emph{recurrent} if $F=1$ and \emph{transient} if $F<1$.

\begin{theorem}[recurrence criterion]\label{thm:recurrence}
$P(V\ge k)=F^k$ for every $k\ge0$, and $\mathbb E[V]=\sum_{n\ge1}u_n$. Consequently:
\begin{enumerate}
\item[(a)] if $\sum_nu_n=\infty$, then $F=1$ and the walk returns to $0$ infinitely often, almost surely;
\item[(b)] if $\sum_nu_n<\infty$, then $F<1$, the walk returns only finitely often, almost surely, and $\sum_{n\ge1}u_n=\frac F{1-F}$.
\end{enumerate}
\end{theorem}
\begin{proof}
Let $T_k$ be the time of the $k$-th return ($T_k=\infty$ if there are fewer than $k$), and $f_n=P(T=n)$. We show by induction that $P(T_k=n)=f^{*k}_n$,
the $k$-fold convolution ($f^{*1}=f$ and $f^{*(k+1)}_n=\sum_mf^{*k}_mf_{n-m}$). The event $\{T_{k+1}=n\}$ is the disjoint union over $m<n$ of
$\{T_k=m\}\cap\{$the walk $(S_{m+j}-S_m)_{j\ge0}$ first returns to $0$ at time $n-m\}$. The first event depends on $X_1,\dots,X_m$ and the second on
$X_{m+1},\dots,X_n$, with probability $f_{n-m}$; so $P(T_{k+1}=n)=\sum_mf^{*k}_mf_{n-m}$. Summing over $n$ (a Cauchy product of non-negative series),
\[
\begin{aligned}P(V\ge k)&=P(T_k<\infty)=\sum_nf^{*k}_n\eqbr=\Bigl(\sum_nf_n\Bigr)^k=F^k.\end{aligned}
\]
Then $\mathbb E[V]=\sum_{k\ge1}P(V\ge k)=\sum_{k\ge1}F^k$, which is $\frac F{1-F}$ if $F<1$ and $\infty$ if $F=1$; and also $\mathbb E[V]=\sum_{n\ge1}\mathbb E[\mathbf 1_{S_n=0}]=\sum_{n\ge1}u_n$ by
monotone convergence. So $\sum u_n=\infty$ if and only if $F=1$, and then $P(V=\infty)=\lim_kP(V\ge k)=1$. If $F<1$, $P(V=\infty)=\lim_kF^k=0$.
\end{proof}

\begin{theorem}[Pólya, 1921]\label{thm:polya}
The simple random walk on $\Z^d$ is recurrent for $d=1$ and $d=2$, and transient for $d\ge3$.
\end{theorem}
\begin{proof}
\emph{$d=1$.} By Lemma~\ref{lem:central}, $\sum_nu_{2n}\ge\sum_n\frac1{2\sqrt n}=\infty$.

\emph{$d=2$.} Turn the plane by $45^\circ$: map $(x,y)$ to $(x+y,\,x-y)$, a bijection of $\Z^2$ onto its image that sends $0$ to $0$ and the steps
$e_1,-e_1,e_2,-e_2$ to $(1,1),(-1,-1),(1,-1),(-1,1)$. So the image of $X_i$ is $(\eta_i,\eta'_i)$, uniform on $\{\pm1\}^2$: the signs $\eta_i$ and $\eta'_i$ are
independent fair signs. The image of $S_n$ is a pair of \emph{independent} one-dimensional walks, and
\[
P(S_{2n}=0)=u_{2n}^2\ge\frac1{4n}
\]
by Lemma~\ref{lem:central}, and the sum of these diverges.

\emph{$d\ge3$.} Realise the steps as $X_i=\varepsilon_ie_{A_i}$, with $A_i$ uniform on $\{1,\dots,d\}$ and $\varepsilon_i$ a fair sign, all independent; this is
uniform on the $2d$ neighbours. Fix $n$ and let $N_a$ be the number of $i\le2n$ with $A_i=a$. Given $(A_1,\dots,A_{2n})=\alpha$, the $a$-th
coordinate of $S_{2n}$ is a sum of $N_a$ independent fair signs, over disjoint sets of indices for different $a$. So
\[
\begin{aligned}P\bigl(S_{2n}=0\mid(A_i)_{i\le2n}=\alpha\bigr)&=\prod_{a=1}^dP(S^{(1)}_{N_a}=0)\\&\le\prod_{a=1}^d\frac1{\sqrt{N_a+1}},\end{aligned}
\]
where $S^{(1)}$ is a one-dimensional walk and the bound is Lemma~\ref{lem:central}. On the event $G$ that $N_a\ge\frac nd$ for every $a$, the product
is at most $(\frac dn)^{d/2}$; off $G$ it is at most $1$. Each $N_a$ is binomial$(2n,\frac1d)$ with mean $\frac{2n}d$, so by Hoeffding's inequality
(Theorem~\ref{thm:hoeffding}), $P(N_a<\frac nd)\le\exp\bigl(-\frac{2(n/d)^2}{2n}\bigr)=e^{-n/d^2}$. Averaging over $\alpha$,
\[
u_{2n}\le\Bigl(\frac dn\Bigr)^{d/2}+d\,e^{-n/d^2},
\]
which is summable when $d/2>1$. As $u_n=0$ for odd $n$, $\sum u_n<\infty$ and Theorem~\ref{thm:recurrence}(b) applies.
\end{proof}

\begin{figure}[ht]
\centering
\fig{ch07-walks}
\caption{$5000$ steps of the walk on $\Z^2$ (left) and on $\Z^3$, projected on a plane (right), from the black dot to the circle. At
this scale they look alike; the difference is in the returns. The planar walk comes back to its start infinitely often,
about $\frac1\pi\ln n$ times by time $n$ (here twice), while the walk in space returns only finitely often, $0.52$ times on average.}
\label{fig:walks}
\end{figure}

\begin{remark}
In dimension $3$ the return probability is $F\approx0.3405$, computed by Watson (1939) and by McCrea and Whipple (1940)
from an integral of the characteristic function: by Fourier inversion on $\Z^d$ (Tome~4), $u_n=\frac1{(2\pi)^d}\int_{[-\pi,\pi]^d}\varphi(t)^n\dd t$
with $\varphi(t)=\frac1d\sum_a\cos t_a$, and summing the geometric series gives $\sum_nu_n=\frac1{(2\pi)^d}\int\frac{\dd t}{1-\varphi(t)}$, which is finite exactly
when $d\ge3$ because $1-\varphi(t)\approx\frac{|t|^2}{2d}$ near $0$. This is the proof that generalises to all walks on $\Z^d$ (Tome~10).
\end{remark}

\begin{beauty}
Recurrence is a property of a series: the walk comes back forever if and only if $\sum_nP(S_n=0)$ diverges, and in $\Z^d$
these probabilities are about $n^{-d/2}$. The dimension enters only through this exponent, and $\sum n^{-d/2}$ converges exactly
when $d\ge3$. A drunkard on a city grid always finds his door again; a drunk bird in the sky may not.
\end{beauty}

\section{Harmonic functions and the Dirichlet problem}

Let $D$ be a finite non-empty subset of $\Z^d$, and $\partial D$ its \emph{boundary}: the points outside $D$ with a neighbour in $D$. A
function $f$ on $\bar D=D\cup\partial D$ is \emph{harmonic in $D$} if, at every $x\in D$, it equals the average of its values at the $2d$ neighbours:
\[
f(x)=\frac1{2d}\sum_{a=1}^d\bigl(f(x+e_a)+f(x-e_a)\bigr).
\]
This is the discrete version of $\Delta f=0$. The \emph{Dirichlet problem} asks for a harmonic $f$ in $D$ with prescribed values $g$ on $\partial D$:
think of a metal plate whose edge is held at temperatures $g$, and of $f$ as the equilibrium temperature inside.

\begin{theorem}[discrete Dirichlet problem]\label{thm:dirichlet}
For every $g:\partial D\to\R$ there is exactly one function $f$ on $\bar D$, harmonic in $D$, with $f=g$ on $\partial D$. It is
\[
\begin{gathered}f(x)=\mathbb E\bigl[g(x+S_\tau)\bigr],\qq\tau=\min\{n\ge0:x+S_n\notin D\},\end{gathered}
\]
the average value of $g$ at the point where the walk from $x$ first leaves $D$. Moreover $\min g\le f\le\max g$.
\end{theorem}
\begin{proof}
\emph{Maximum principle.} Let $f$ be harmonic in $D$ and let $M$ be its maximum over the finite set $\bar D$, attained at $x_0$. If $x_0\in D$,
then $M=f(x_0)$ is the average of $2d$ values that are all $\le M$, so they all equal $M$; in particular $f(x_0+e_1)=M$. Repeating, $f=M$ at
$x_0+ke_1$ as long as $x_0+(k-1)e_1\in D$; since $D$ is finite, some $x_0+ke_1$ lies in $\partial D$. So the maximum is attained on $\partial D$, and in the same way
the minimum. Hence $\min_{\partial D}f\le f\le\max_{\partial D}f$.

\emph{Uniqueness.} The difference of two solutions is harmonic in $D$ and zero on $\partial D$, so its maximum and minimum are $0$.

\emph{Existence.} The first coordinates of the points of $D$ lie in an interval of length $L$, so $L+1$ consecutive steps $+e_1$ take the walk
out of $D$, from any starting point in $D$. As in Lemma~\ref{lem:exit}, with blocks of $L+1$ steps, $P(\tau>k(L+1))\le\bigl(1-(2d)^{-(L+1)}\bigr)^k$, so $\tau<\infty$
almost surely and $x+S_\tau\in\partial D$; $f$ is well defined, as $g$ is bounded. If $x\in\partial D$ then $\tau=0$ and $f(x)=g(x)$. If $x\in D$ then $\tau\ge1$, and
conditioning on the first step exactly as in the proof of Theorem~\ref{thm:ruin} gives $f(x)=\frac1{2d}\sum_{y\sim x}f(y)$: after the first step, the
walk is a fresh walk started from the neighbour $y=x+X_1$, independent of $X_1$.
\end{proof}

The gambler's ruin with $p=\frac12$ is the case $d=1$, $D=\{1,\dots,N-1\}$, $g(0)=0$, $g(N)=1$: the harmonic functions on an interval are
the affine ones, and $h(x)=\frac xN$.

\begin{figure}[ht]
\centering
\fig{ch07-dirichlet}
\caption{The solution of the Dirichlet problem on an L-shaped domain of $\Z^2$, with $g=1$ on the top edge (orange) and
$g=0$ on the rest of the boundary. The shade at $x$ is $f(x)$, the probability that a walk from $x$ leaves through the top
edge. The two black paths are walks from the marked point, one leaving through the top and one through the side.}
\label{fig:dirichlet}
\end{figure}

\begin{remark}
Refine the grid to mesh $\frac1k$ and let $k\to\infty$: the walk becomes Brownian motion (Chapter~11) and the discrete harmonic
functions become solutions of Laplace's equation $\Delta f=0$. Kakutani showed in 1944 that the solution of the continuous Dirichlet
problem is $f(x)=\mathbb E[g(B_\tau)]$, with $B$ a Brownian motion from $x$ and $\tau$ its exit time: Theorem~\ref{thm:dirichlet}, word for word. The
proof needs stochastic calculus and is in Tome~11. Harmonic functions on graphs are also the voltages of electrical
networks, which gives a second, physical view of recurrence (Exercise~\ref{ex:electric}).
\end{remark}

\section{A random walk that solves 2-SAT}

A \emph{2-SAT formula} in boolean variables $v_1,\dots,v_n$ is a conjunction of clauses $(\ell\vee\ell')$, where each literal $\ell,\ell'$ is a variable or
its negation, for example $(v_1\vee\neg v_2)\wedge(v_2\vee v_3)\wedge(\neg v_1\vee\neg v_3)$. Deciding whether such a formula can be satisfied is the
simplest interesting case of the satisfiability problem. Christos Papadimitriou found in 1991 an algorithm that
is nothing but a random walk:
\begin{quote}
Start from any assignment. While some clause is false, pick a false clause (by any rule at all), choose one of
its two literals uniformly at random, and flip the value of its variable.
\end{quote}

\begin{theorem}[Papadimitriou]\label{thm:twosat}
If the formula is satisfiable, the expected number of flips before the algorithm finds a satisfying assignment is
at most $n^2$, whatever the starting assignment and whatever the rule for picking false clauses.
\end{theorem}
\begin{proof}
Fix a satisfying assignment $A^*$, let $Y_t$ be the assignment after $t$ flips, $X_t$ the number of variables on which $Y_t$ agrees with $A^*$,
and $T$ the first $t$ at which $Y_t$ satisfies the formula. Since $X_t=n$ means $Y_t=A^*$, $T\le\min\{t:X_t=n\}$.

Each flip changes one variable, so $X_{t+1}=X_t\pm1$. If $T>t$, the chosen clause is false under $Y_t$, so both its literals are false; $A^*$
makes at least one of them true, so at least one of the two variables disagrees with $A^*$, and flipping it increases $X$. So,
conditionally on any history of the first $t$ flips with $T>t$ and $X_t=x$, the probability $\theta$ that $X_{t+1}=x+1$ is at least $\frac12$ (and $\theta=1$
if $x=0$).

Let $\Psi(x)=n^2-x^2$. On such a history,
\[
\begin{aligned}&\mathbb E[\Psi(X_{t+1})+1\mid\text{history}]-\Psi(x)\\&\qquad=-\theta(2x+1)+(1-\theta)(2x-1)+1\\&\qquad=2x(1-2\theta)\le0.\end{aligned}
\]
Let $Z_t=\Psi(X_{t\wedge T})+t\wedge T$. Summing over the finitely many histories of the first $t$ flips (on those with $T\le t$, $Z_{t+1}=Z_t$), we get
\[
\mathbb E[Z_{t+1}]\le\mathbb E[Z_t]\le\dots\le\mathbb E[Z_0]=\Psi(X_0).
\]
As $\Psi\ge0$ on $[0,n]$, $\mathbb E[t\wedge T]\le\Psi(X_0)\le n^2$ for every $t$, and monotone convergence gives $\mathbb E[T]\le n^2$.
\end{proof}

\begin{corollary}\label{cor:twosat}
Run the algorithm for $2kn^2$ flips. If the formula is satisfiable, the probability that no satisfying assignment has been
found is at most $2^{-k}$; if it is not, none is ever found.
\end{corollary}
\begin{proof}
By Markov's inequality, from any assignment the probability of needing more than $2n^2$ further flips is at most
$\frac{n^2}{2n^2}=\frac12$. Theorem~\ref{thm:twosat} applies afresh after each block of $2n^2$ flips, from whatever assignment the block ends with, and
conditionally on the past (the proof used only the rule for the next flip). So each of the $k$ blocks fails with
conditional probability at most $\frac12$.
\end{proof}

The proof is a comparison with the walk on $\{0,\dots,n\}$ that goes up with probability $\frac12$ and is reflected at $0$, whose
expected time to reach $n$ from $x$ is exactly $n^2-x^2$. The bound is sharp: for the formula made of $(v_1\vee v_1)$ and the clauses
$(\neg v_i\vee v_{i+1})$ and $(v_i\vee\neg v_{i+1})$ for $i<n$, started from all variables false, the assignment is always ``$v_1,\dots,v_k$ true, the
others false'' with a single false clause, and $k$ performs exactly that reflected walk. For 3-SAT the same argument only gives $\theta\ge\frac13$, a walk drifting the
wrong way, and an exponential time; Schöning (1999) turned this into the algorithm with restarts that solves 3-SAT
in time about $(\frac43)^n$, still one of the best known.

\section{Exercises}

\begin{exercise}[$\star$]\label{ex:ballotwalk}
Show that $P(S_1>0,\dots,S_n>0\mid S_n=k)=\frac kn$ for $0<k\le n$ with $n+k$ even.
\end{exercise}

\begin{exercise}[$\star\star$]\label{ex:maxlaw}
Show that $\frac{M_n}{\sqrt n}\to|N|$ in law, where $N$ is $N(0,1)$ and $M_n$ is the maximum of Proposition~\ref{prop:maximum}: that is,
$P(M_n\ge t\sqrt n)\to2(1-\Phi(t))$ for every $t>0$.
\end{exercise}

\begin{exercise}[$\star\star\star$]\label{ex:timepositive}
\emph{The other arcsine law.} Say that the $k$-th step of the walk is \emph{positive} if $S_{k-1}+S_k>0$ (the segment from $S_{k-1}$ to $S_k$
lies above the axis). Let $P_{2n}$ be the number of positive steps among the first $2n$. Show that $P(P_{2n}=2k)=u_{2k}u_{2n-2k}$ for
$0\le k\le n$ (Chung and Feller, 1949).
\end{exercise}

\begin{exercise}[$\star$]\label{ex:unfair}
In the gambler's ruin with $p<\frac12$ and even $x$ and $N$, compare the probability of reaching $N$ from $x$ with stakes of
$1$ euro and of $2$ euros. Which is larger, and by what factor?
\end{exercise}

\begin{exercise}[$\star\star$]\label{ex:meet}
Two independent simple random walks on $\Z^d$ start at $0$. Show that the probability that they are at the same
point at time $n$ is $u_{2n}$, and deduce that the expected number of meetings is infinite if and only if $d\le2$.
\end{exercise}

\begin{exercise}[$\star\star$]\label{ex:escape}
Let $d\ge3$. Show that $P(S_n=x)\le u_{2\lfloor n/2\rfloor}$ for every $x\in\Z^d$ and every $n$, and deduce that $|S_n|\to\infty$ almost surely: the walk
visits each point only finitely often.
\end{exercise}

\begin{exercise}[$\star\star$]\label{ex:exittime}
Let $\tau$ be the first time the walk on $\Z^d$ started at $0$ reaches a point with $|S_n|\ge R$ (Euclidean norm, $R>0$). Show that
$R^2\le\mathbb E[\tau]<(R+1)^2$, in every dimension.
\end{exercise}

\begin{exercise}[$\star\star$]\label{ex:electric}
\emph{Electrical networks.} Let $G$ be a finite connected graph and $a\ne b$ two vertices. The walk on $G$ moves to a uniform
neighbour at each step. Let $h(x)$ be the probability that the walk from $x$ reaches $b$ before $a$.
\begin{enumerate}
\item[(a)] Show that $h$ is the unique function with $h(a)=0$, $h(b)=1$ that equals the average of its values at the
neighbours at every other vertex. (With unit resistors on the edges, $h$ is the voltage when $a$ is earthed and $b$ held
at $1$ volt.)
\item[(b)] Let $I=\sum_{y\sim a}h(y)$ (the current leaving $b$ and entering $a$) and $R=\frac1I$, the effective resistance. Show that the walk from $a$
reaches $b$ before returning to $a$ with probability $\frac1{\deg(a)R}$.
\item[(c)] Compute $R$ when $G$ is a path of length $m$ from $a$ to $b$, and when $G$ is a cycle in which $a$ and $b$ are joined by two
paths of lengths $m$ and $k$. Recognise the rules for resistances in series and in parallel.
\end{enumerate}
\end{exercise}

\section{Hints}

\begin{hint}{ex:ballotwalk}
Given $S_n=k$, all orders of the $\frac{n+k}2$ up-steps and $\frac{n-k}2$ down-steps are equally likely: this is the ballot theorem.
\end{hint}
\begin{hint}{ex:maxlaw}
Use Proposition~\ref{prop:maximum} with $r=\lceil t\sqrt n\rceil$, the central limit theorem, and Lemma~\ref{lem:central} for the term $P(S_n=r)$.
\end{hint}
\begin{hint}{ex:timepositive}
Induction on $n$. The cases $k=0$ and $k=n$ are Theorem~\ref{thm:mainlemma}(b). Otherwise the walk returns to $0$ before time $2n$; split according
to the first return time $2r$ and the sign of the first excursion, and use Lemma~\ref{lem:renewal}.
\end{hint}
\begin{hint}{ex:unfair}
With stakes of $2$, the game goes from $\frac x2$ to $\frac N2$. Factor $r^x-1=(r^{x/2}-1)(r^{x/2}+1)$.
\end{hint}
\begin{hint}{ex:meet}
$P(S_n=S'_n)=\sum_yP(S_n=y)P(S'_n=y)$, and $-S'_n$ has the same law as $S'_n$.
\end{hint}
\begin{hint}{ex:escape}
For $n=2m$, split $S_{2m}=S_m+(S_{2m}-S_m)$ and use the Cauchy--Schwarz inequality together with Exercise~\ref{ex:meet}. Then use the first
Borel--Cantelli lemma.
\end{hint}
\begin{hint}{ex:exittime}
$\mathbb E\bigl[|S_{n+1}|^2-|S_n|^2\mid X_1,\dots,X_n\bigr]=1$. Show that $\mathbb E\bigl[|S_{t\wedge\tau}|^2\bigr]=\mathbb E[t\wedge\tau]$ for every $t$, then let $t\to\infty$.
\end{hint}
\begin{hint}{ex:electric}
(a) Copy the proof of Theorem~\ref{thm:dirichlet}, with $D$ the set of vertices other than $a$ and $b$. (b) First-step analysis from $a$.
(c) On a path, the harmonic functions are affine.
\end{hint}

\section{Solutions}

\begin{solution}{ex:ballotwalk}
\idea Conditioning on the endpoint makes all paths with that endpoint equally likely, and that is the setting of the
ballot theorem.
\writeup Each path of length $n$ has probability $2^{-n}$, so given $S_n=k$ the walk is uniform among the paths with $a=\frac{n+k}2$ up-steps
and $b=\frac{n-k}2$ down-steps, that is, a uniform order of $a$ votes for A and $b$ for B. By the ballot theorem (Theorem~\ref{thm:ballot}),
the probability that A stays strictly ahead is $\frac{a-b}{a+b}=\frac kn$.
\end{solution}

\begin{solution}{ex:maxlaw}
\idea The maximum is ``twice the tail'' of $S_n$ by reflection, and the tail of $S_n$ is Gaussian by the CLT.
\writeup Fix $t>0$ and let $r=\lceil t\sqrt n\rceil$, so $\{M_n\ge t\sqrt n\}=\{M_n\ge r\}$ and, by Proposition~\ref{prop:maximum},
\[
P(M_n\ge t\sqrt n)=2P(S_n>r)+P(S_n=r).
\]
By Lemma~\ref{lem:central}, $P(S_n=r)\le\frac1{\sqrt{n+1}}\to0$. Let $\varepsilon>0$; for $n>\varepsilon^{-2}$, $t\sqrt n\le r<(t+\varepsilon)\sqrt n$, so
\[
\begin{aligned}P\Bigl(\frac{S_n}{\sqrt n}>t+\varepsilon\Bigr)&\le P(S_n>r)\eqbr\le P\Bigl(\frac{S_n}{\sqrt n}\ge t\Bigr)\eqbr\le P\Bigl(\frac{S_n}{\sqrt n}>t-\varepsilon\Bigr).\end{aligned}
\]
By the central limit theorem (Corollary~\ref{cor:clt}, with $\sigma=1$) the two outer terms tend to $1-\Phi(t+\varepsilon)$ and $1-\Phi(t-\varepsilon)$. Letting $\varepsilon\to0$,
$P(S_n>r)\to1-\Phi(t)$ and $P(M_n\ge t\sqrt n)\to2(1-\Phi(t))=P(|N|\ge t)$. As the limit is continuous in $t$, this gives the convergence of the
distribution functions at every point (squeeze $P(\frac{M_n}{\sqrt n}<t)$ between $1-P(\frac{M_n}{\sqrt n}\ge t\mp\varepsilon)$), which is convergence in law.
\end{solution}

\begin{solution}{ex:timepositive}
\idea Before its first return the walk is entirely on one side, and after it the walk starts afresh: this gives a
recursion that the product $u_{2k}u_{2n-2k}$ satisfies, thanks to the first-return decomposition.
\writeup Between two consecutive zeros the walk stays on one side, so the first $2n$ steps split into pieces of even length (the
times of zeros are even, and so is $2n$), each made only of positive or only of negative steps. Hence $P_{2n}$ is even; let
$b_{2k,2n}=P(P_{2n}=2k)$. We prove $b_{2k,2n}=u_{2k}u_{2n-2k}$ by induction on $n$, with $b_{j,m}=0$ for $j<0$ or $j>m$.

If $k=n$, all steps are positive exactly when $S_1\ge0,\dots,S_{2n}\ge0$, of probability $u_{2n}$ by Theorem~\ref{thm:mainlemma}(b); $k=0$ is symmetric.
This covers $n=1$. Let $n\ge2$ and $1\le k\le n-1$. Then the walk has both positive and negative steps, so it changes side and
returns to $0$ at some time $T=2r\le2n$. Given $T=2r$, the first $2r$ steps are all positive or all negative, each with probability
$\frac12$ (reflecting the first $2r$ steps preserves $\{T=2r\}$), and the remaining $2n-2r$ steps form an independent walk. Hence
\[
\begin{aligned}b_{2k,2n}&=\frac12\sum_{r=1}^nf_{2r}\,b_{2k-2r,2n-2r}\eqbr+\frac12\sum_{r=1}^nf_{2r}\,b_{2k,2n-2r}.\end{aligned}
\]
By the induction hypothesis and Lemma~\ref{lem:renewal} (which applies as $k\ge1$), the first sum is
\[
\sum_{r=1}^kf_{2r}\,u_{2k-2r}\,u_{2n-2k}=u_{2k}\,u_{2n-2k},
\]
and in the same way (as $n-k\ge1$) the second is
\[
\sum_{r=1}^{n-k}f_{2r}\,u_{2k}\,u_{2n-2k-2r}=u_{2k}\,u_{2n-2k}.
\] So $b_{2k,2n}=u_{2k}u_{2n-2k}$.
\end{solution}

\begin{solution}{ex:unfair}
\idea Doubling the stake halves the number of bets the house gets to take its edge from.
\writeup Let $r=\frac qp>1$. With stakes of $1$, the success probability is $\frac{r^x-1}{r^N-1}$; with stakes of $2$, it is the same game from $\frac x2$ to
$\frac N2$, so $\frac{r^{x/2}-1}{r^{N/2}-1}$. Their quotient is
\[
\frac{r^x-1}{r^N-1}\cdot\frac{r^{N/2}-1}{r^{x/2}-1}=\frac{r^{x/2}+1}{r^{N/2}+1}<1.
\]
So bold play wins: the larger stake is better, by the factor $\frac{r^{N/2}+1}{r^{x/2}+1}$, about $r^{(N-x)/2}$. For the roulette example ($x=100$, $N=200$) the
factor is about $14$. (For $p>\frac12$ the inequality reverses, and for $p=\frac12$ both are $\frac xN$.)
\end{solution}

\begin{solution}{ex:meet}
\idea The difference of two independent walks is a walk run at double speed.
\writeup By independence and symmetry,
\[
\begin{aligned}P(S_n=S'_n)&=\sum_yP(S_n=y)P(S'_n=y)\\&=\sum_yP(S_n=y)P(-S'_n=-y)\\&=P(S_n-S'_n=0).\end{aligned}
\]
Now $-S'_n$ is a sum of $n$ independent uniform steps, independent of $S_n$, so $S_n+(-S'_n)$ has the law of $S_{2n}$, and the probability is $u_{2n}$.
The expected number of meetings at times $n\ge1$ is $\sum_{n\ge1}u_{2n}$, which is $\sum_{n\ge1}u_n$ (odd terms vanish); it is infinite exactly when $d\le2$, by
the proof of Pólya's theorem.
\end{solution}

\begin{solution}{ex:escape}
\idea The most likely place to be at an even time is the starting point; then summability of $u_{2n}$ plus Borel--Cantelli does the rest.
\writeup For $n=2m$, $S_{2m}=S_m+(S_{2m}-S_m)$ with independent terms of the same law, so by Cauchy--Schwarz
\[
\begin{aligned}P(S_{2m}=x)&=\sum_yP(S_m=y)P(S_m=x-y)\\&\le\Bigl(\sum_yP(S_m=y)^2\Bigr)^{1/2}\eqbr\cdot\Bigl(\sum_yP(S_m=x-y)^2\Bigr)^{1/2},\end{aligned}
\]
and both factors equal $\bigl(\sum_yP(S_m=y)^2\bigr)^{1/2}$, while $\sum_yP(S_m=y)^2=u_{2m}$ by Exercise~\ref{ex:meet}. For $n=2m+1$, $P(S_{2m+1}=x)$ is the average of the
$P(S_{2m}=x-e)$ over the $2d$ unit vectors $e$, so it is at most $u_{2m}$ too. Hence $\sum_nP(S_n=x)\le2\sum_mu_{2m}<\infty$ for $d\ge3$, and by the first
Borel--Cantelli lemma (Theorem~\ref{thm:bc}) the walk visits $x$ only finitely often, almost surely. The ball $\{|x|\le R\}$ contains finitely many
points, so almost surely the walk eventually leaves it for good; intersecting over $R=1,2,\ldots$ gives $|S_n|\to\infty$ almost surely.
\end{solution}

\begin{solution}{ex:exittime}
\idea $|S_n|^2-n$ is a ``fair game'', so the expected time to reach distance $R$ is the expected value of $|S_\tau|^2$, which is about $R^2$.
\writeup $|S_{n+1}|^2=|S_n|^2+2S_n\cdot X_{n+1}+1$, since $|X_{n+1}|=1$. The event $\{\tau>n\}$ and $S_n$ depend only on $X_1,\dots,X_n$, and $\mathbb E[X_{n+1}]=0$, so
with $\Delta_n=|S_{(n+1)\wedge\tau}|^2-|S_{n\wedge\tau}|^2$,
\[
\begin{aligned}\mathbb E[\Delta_n]&=\mathbb E\bigl[\mathbf 1_{\tau>n}(2S_n\cdot X_{n+1}+1)\bigr]\eqbr=P(\tau>n),\end{aligned}
\]
which is also $\mathbb E[(n+1)\wedge\tau-n\wedge\tau]$. Summing, $\mathbb E\bigl[|S_{t\wedge\tau}|^2\bigr]=\mathbb E[t\wedge\tau]$ for every $t$. The walk exits the finite set $\{|x|<R\}$ in finite time
(proof of Theorem~\ref{thm:dirichlet}), and $|S_{t\wedge\tau}|<R+1$ since the step before $\tau$ is inside the ball. Letting $t\to\infty$, dominated convergence on
the left and monotone convergence on the right give $\mathbb E[\tau]=\mathbb E\bigl[|S_\tau|^2\bigr]$, which lies in $[R^2,(R+1)^2)$.
\end{solution}

\begin{solution}{ex:electric}
\idea Harmonic functions, walks and voltages are three faces of one object: averaging over neighbours is Kirchhoff's
current law with unit resistors.
\writeup (a) The proof of Theorem~\ref{thm:dirichlet} works on any finite connected graph, with $D$ the vertices other than $a$, $b$ and
$\partial D=\{a,b\}$. Maximum principle: from a maximum at $x\in D$, all neighbours are also maximal; follow a path from $x$ to $b$ (the graph is
connected) until it first leaves $D$. Existence: from any vertex, a fixed shortest path to $\{a,b\}$, of length at most $|G|$, is followed
with probability at least $\Delta^{-|G|}$ ($\Delta$ the maximum degree), so the walk reaches $\{a,b\}$ in finite time; first-step analysis shows that $h$ is
harmonic in $D$. (Kirchhoff: the current from $x$ to $y$ is $h(x)-h(y)$, and harmonicity says the currents out of $x$ sum to $0$.)

(b) From $a$, the walk moves to a uniform neighbour $y$ and then, as a fresh walk, reaches $b$ before $a$ with probability $h(y)$. So the
probability is $\frac1{\deg a}\sum_{y\sim a}h(y)=\frac I{\deg a}=\frac1{\deg(a)R}$.

(c) On a path $a=x_0,\dots,x_m=b$, $h(x_j)=\frac jm$ is harmonic, so $I=\frac1m$ and $R=m$: resistances in series add. On the cycle, $h$ is affine on each of
the two paths, so the two neighbours of $a$ have $h=\frac1m$ and $\frac1k$, giving $I=\frac1m+\frac1k$ and $R=\frac{mk}{m+k}$: conductances in parallel add. For a path
from $a$, (b) gives escape probability $\frac1m$, the gambler's ruin again.
\end{solution}

\section*{Sources and further reading}

\begin{itemize}
\item W.~Feller, \emph{An Introduction to Probability Theory and Its Applications}, Vol.~1, 3rd ed., Wiley, 1968,
Chapters~III and~XIV. The main lemma, the arcsine laws and the gambler's ruin, with the path-counting style of this
chapter.
\item G.~Pólya, \emph{Über eine Aufgabe der Wahrscheinlichkeitsrechnung betreffend die Irrfahrt im Straßennetz},
Mathematische Annalen 84 (1921), 149--160. The recurrence theorem.
\item K.~L.~Chung and W.~Feller, \emph{On fluctuations in coin-tossing}, Proceedings of the National Academy of Sciences 35
(1949), 605--608. The arcsine law for the time spent positive.
\item P.~G.~Doyle and J.~L.~Snell, \emph{Random Walks and Electric Networks}, Mathematical Association of America, 1984
(free at arXiv:math/0001057). Recurrence through resistance, a delightful short book.
\item G.~F.~Lawler and V.~Limic, \emph{Random Walk: A Modern Introduction}, Cambridge University Press, 2010. Return
probabilities, harmonic functions and the Dirichlet problem on $\Z^d$ in depth.
\item C.~H.~Papadimitriou, \emph{On selecting a satisfying truth assignment}, Proceedings of the 32nd IEEE Symposium on
Foundations of Computer Science (1991), 163--169. The 2-SAT walk.
\item U.~Schöning, \emph{A probabilistic algorithm for $k$-SAT and constraint satisfaction problems}, Proceedings of the
40th IEEE Symposium on Foundations of Computer Science (1999), 410--414.
\item M.~Mitzenmacher and E.~Upfal, \emph{Probability and Computing}, 2nd ed., Cambridge University Press, 2017,
Section~7.1: 2-SAT and 3-SAT by random walks.
\end{itemize}
